\chapter{Discussion, Conclusion and Future Work}
\label{chap:conclusion}

\section{Discussion of Results}

This project investigated language-guided active view selection for multi-view object perception. The main goal was to test whether CLIP-based scoring, a learned MLP reranker, and a trained active view selector can select views that satisfy a natural language instruction under a limited observation budget.

The results provide clear support for language-guided scoring and stronger support for learned reranking. On both CO3D and UMD+GT, CLIP-only methods outperform non-semantic random selection, but they remain close to exhaustive CLIP. This indicates that image--text similarity is useful, yet cosine similarity alone is not a strong enough final ranker for fine-grained viewpoint instructions.

The learned MLP reranker is the central improvement. By adapting frozen CLIP features to instruction-specific ground-truth views, it raises accuracy substantially as soon as multiple views are observed. On shuffled CO3D, random selection with learned reranking reaches 0.478 at budget~3 and 0.520 at budget~5, compared with about 0.31 for random-plus-CLIP. On UMD+GT, the same pattern is even stronger: learned reranking alone reaches 0.632 at budget~3 and 0.696 at budget~5.

The value of the trained selector depends on the dataset. On UMD+GT, over 20 trials, the full pipeline of trained selection plus learned reranking is clearly best at low budgets, reaching \(0.669\pm0.011\) already at budget~2 and \(0.680\pm0.008\) at budget~3. Affordance visibility changes smoothly around the turntable, so the selector can learn which next view is likely to be useful. On CO3D, the full pipeline helps at budget~2, but CLIP relevance-plus-novelty combined with learned reranking can match or exceed it at budget~3. This is consistent with the fact that the selector and reranker share the same training objective and therefore tend to reinforce the same view preferences rather than explore complementary ones.

Ablation studies make these contributions explicit. Every plotted gain is defined as accuracy(method) minus accuracy(random selection + CLIP rerank). Across both datasets, selection-only methods add only a few overall points, while any configuration that includes the learned reranker produces double-digit gains. The practical implication is that, for the current problem formulation, improving final observed-view ranking is more important than refining the CLIP heuristic alone.

\section{Limitations}

This project has several limitations.

First, the selector and reranker are trained with the same binary ``is this view ground truth?'' objective. They therefore learn overlapping signals. A more natural active-perception objective would train the selector to choose the view that most improves the expected final answer given the views already observed.

Second, candidate views are scored using their true image features before they have been ``visited''. This is acceptable for offline ranking from a pre-collected view set, but it is optimistic for embodied robotics, where the next image is unavailable until the camera moves. A realistic system would need to predict the value of a candidate pose from the current observation and camera geometry.

Third, both datasets remain modest in scale. CO3D evaluation uses 50 samples (10 per category). UMD+GT has 345 test samples overall, but some affordances such as contain and wrap-grasp have only eight test samples after filters remove rotationally ambiguous cases. Per-group differences are therefore noisy.

Fourth, ground-truth definitions are proxies. CO3D labels depend on human judgement of instruction satisfaction. UMD+GT labels use relative affordance mask area, which correlates with visibility but is not identical to human notions of the most informative viewpoint.

Fifth, the system selects from pre-collected views and does not control a physical robot. It evaluates the decision component of active perception, not continuous camera motion, calibration, or closed-loop control.

Finally, prompt design is fixed. Because the system depends on text embeddings, alternative phrasings of the same instruction could change both CLIP scores and learned-model behaviour.

\section{Future Work}

There are several directions for future work.

First, the selector should be trained with a sequential decision objective. Instead of predicting whether a view is ground truth in isolation, it could estimate the expected accuracy gain of adding a view given the current observed set and the fixed reranker.

Second, next-view prediction should be made pose-conditioned rather than image-oracle. A model that scores candidate camera poses without access to their images would better match real robotic constraints.

Third, the datasets should be expanded. More CO3D categories, more UMD tools after less aggressive filtering, and graded relevance labels would make evaluation more robust.

Fourth, stronger multimodal backbones could be tested. Larger vision-language models with better spatial grounding may improve both zero-shot CLIP baselines and the quality of frozen features used by the MLP.

Fifth, the system could be extended to a real robot. The active policy would then choose the next camera pose, capture a new image, update the observed set, and stop when the budget or a confidence criterion is reached.

Finally, richer language supervision could be studied, including multiple prompt templates, affordance-specific prompts, and instructions that combine viewpoint and action goals such as ``find a view suitable for grasping the handle''.

\section{Conclusion}

This thesis presented a language-guided active view selection pipeline for multi-view object perception. The system starts from a random view, selects additional views under a limited budget, and reranks the observed views to produce a final prediction. The pipeline includes a CLIP relevance-plus-novelty policy, a learned MLP reranker, and a trained active view selector, evaluated on a shuffled CO3D instruction dataset and an affordance-derived UMD+GT dataset.

The main conclusion is that language-guided scoring is useful, but a lightweight learned reranker on frozen CLIP features is far more effective than CLIP cosine similarity alone. On UMD+GT, combining the trained selector with the learned reranker further improves low-budget performance, reaching about 0.67--0.68 accuracy by budgets 2--3 against a random-plus-CLIP baseline near 0.46, with stable variance across 20 trials. On CO3D, learned reranking remains the dominant contributor, while novelty-aware selection can complement it as well as or better than a selector trained with the same GT-view objective.

Overall, this project provides a practical experimental foundation for language-guided active perception. With sequential selector training, pose-conditioned next-view prediction, larger datasets, and embodied evaluation, active view selection can become a stronger component of future robotic vision-language systems.
